\subsection{Wavelets: spectral selection with spatial localization}

Fourier modes answer which graph frequencies are present, but each eigenvector is generally spread
over the whole graph.  A wavelet asks a more localized question: at what scale, and around which
node, does a change occur?  On a self-adjoint graph Laplacian
$L=V\Lambda V^{\mathsf T}$, choose a band-pass generating function $\psi$ and define
\begin{equation}
\Psi_s
=\psi(sL)
=V\psi(s\Lambda)V^{\mathsf T},
\qquad
\psi_{s,i}=\Psi_s\delta_i.
\label{eq:graph-wavelet-atom}
\end{equation}
The vector $\psi_{s,i}$ is a graph-wavelet atom centered at node $i$ and scale $s$
\cite{hammond2011}.  Its spectrum is selected by $\psi(s\lambda)$, while its node-space profile
records how that band is expressed around $i$.  Graph wavelets couple spectral selectivity to
spatial localization.

The scale convention depends on the chosen generator.  For the common dilation
$\psi(s\lambda)$, increasing $s$ shifts a fixed band-pass profile toward smaller eigenvalues and
usually produces a broader node-space atom.  Small scales emphasize finer graph variation; large
scales emphasize coarser organization.  This is a statement relative to $L$.  Atomic-motion and
protein-domain interpretations require a separately calibrated correspondence between graph scale
and the biological geometry.

Localization follows from the spectral construction.  If $\psi(s\lambda)$ is approximated by a
degree-$K$ polynomial,
\begin{equation}
\psi(sL)\approx\sum_{k=0}^{K}c_k(s)L^k,
\label{eq:graph-wavelet-polynomial}
\end{equation}
then $\psi(sL)\delta_i$ is supported, up to approximation error, within $K$ graph steps of $i$.
This is the same spectral-to-spatial passage used by Chebyshev graph filters.  The wavelet family
organizes several band-pass scales around localized centers; a task-specific graph filter instead
learns the propagation profile required by its objective.

A low-pass scaling function $\phi_J$ retains the coarsest component.  A stable wavelet transform
requires more than an appealing collection of bands.  If constants $0<A\leq B<\infty$ satisfy
\begin{equation}
A
\leq
|\phi_J(\lambda)|^2+\sum_{s\in\mathcal S}|\psi(s\lambda)|^2
\leq B
\label{eq:graph-wavelet-frame-spectrum}
\end{equation}
on the spectrum of $L$, then
\begin{equation}
A\|y\|_2^2
\leq
\|\phi_J(L)y\|_2^2
+\sum_{s\in\mathcal S}\|\psi(sL)y\|_2^2
\leq B\|y\|_2^2.
\label{eq:graph-wavelet-frame}
\end{equation}
The transform is then a stable frame: it can be redundant without losing control of the signal
energy.  This distinguishes multiscale analysis from an arbitrary collection of spectral plots.

Heat diffusion, wavelets, and graph convolution can now be placed in one functional calculus:
\begin{equation}
\begin{aligned}
e^{-tL}&:&\quad&\text{low-pass smoothing at diffusion scale }t,\\
\psi(sL)&:&\quad&\text{localized band-pass analysis at scale }s,\\
\sum_{k=0}^{K}\theta_kL^k&:&\quad&\text{finite-hop polynomial filtering},\\
S^KH&:&\quad&\text{repeated propagation with possible over-smoothing}.
\end{aligned}
\label{eq:spectral-operator-family}
\end{equation}
Their common mathematics is $g(L)$; their state, objective, and interpretation remain different.
A heat kernel describes smoothing relative to a graph, a wavelet transform analyzes a signal, and a
GCN layer is trained as part of a predictor.

For an operator-valued protein graph, the analyzed signal must still be declared.  It might be a
scalar conservation score, a static interaction norm, a mutational response field, or a categorical
section $f\in\bigoplus_i\mathcal H_i^0$.  Scalar signals can be analyzed with a scalar Laplacian.
Categorical sections require the connection Laplacian so that neighboring contrasts are transported
between fibers before they are compared:
\begin{equation}
f\longmapsto\psi(s\mathbb L_{\mathrm{conn}})f.
\label{eq:connection-wavelet}
\end{equation}
The resulting coefficient measures scale-dependent coherence in the declared transport geometry.
Recovery of the original edge operator would require an additional invertible representation and
its associated reconstruction map.

Classical spectral graph wavelets also assume a self-adjoint operator with real spectral calculus.
A directed attention route $P$ is generally nonnormal and therefore calls for a different spectral
construction.  One may analyze the reversible
Laplacian, a singular-value construction, or a declared symmetrization, but each discards or
reorganizes directed circulation differently.  Thus the order of reasoning remains
\begin{equation}
\boxed{
\begin{gathered}
\text{declare the graph and signal}
\longrightarrow
\text{choose the operator geometry}\\
\longrightarrow
\text{choose the scale filters}
\longrightarrow
\text{interpret localized coefficients}.
\end{gathered}}
\label{eq:wavelet-audit-order}
\end{equation}
Low-, band-, and high-frequency coefficients are mathematical facts about that construction.
Calling them conservation, allostery, physical vibration, or mechanism requires an external
endpoint that licenses the name.
